
\subsubsection{RQ1: Existence of Task-Category Clustering}

\begin{table}[htbp]
\centering
\resizebox{\columnwidth}{!}{%
\begin{tabular}{ll}
\toprule
Metric & Value \\
\midrule
Within-category mean distance & 7.606 \\
Between-category mean distance & 7.795 \\
Mean difference & 0.190 \\
Cohen's d & 0.765 \\
95\% CI for Cohen's d & [0.68, 0.85] \\
p-value (Mann-Whitney U) & 8.63 $\times 10^{-28}$ \\
Within-category pairs & 380 \\
Between-category pairs & 400 \\
\bottomrule
\end{tabular}
}
\end{table}

The within-category Grassmann distances are smaller than between-category distances. The effect size (Cohen's d = 0.765) exceeds the threshold of 0.5. The 95\% confidence interval for the mean difference [0.155, 0.226] does not include zero.

\begin{figure}[htbp]
\centering
\includegraphics[width=\columnwidth]{figures/distance_heatmap.png}
\caption{Distance Heatmap}
\label{fig:distance_heatmap}
\end{figure}
\textit{Figure 1: Pairwise Grassmann distance matrix between all 40 LoRA adapters.}

\subsubsection{RQ2: Correlation with Semantic Similarity}

\begin{table}[htbp]
\centering
\resizebox{\columnwidth}{!}{%
\begin{tabular}{ll}
\toprule
Metric & Value \\
\midrule
Spearman $\rho$ & 0.389 \\
95\% CI & [0.328, 0.450] \\
p-value & 1.29 $\times 10^{-29}$ \\
N pairs & 780 \\
\bottomrule
\end{tabular}
}
\end{table}

The positive correlation indicates that adapters for tasks in the same FLAN category have smaller Grassmann distances than adapters for tasks in different categories. The 95\% confidence interval excludes zero.

\begin{figure}[htbp]
\centering
\includegraphics[width=\columnwidth]{figures/scatter_regression.png}
\caption{Scatter Regression}
\label{fig:scatter_regression}
\end{figure}
\textit{Figure 2: Grassmann distance vs. FLAN taxonomy distance showing positive correlation ($\rho$ = 0.389).}

\subsubsection{RQ3: Layer-Type Analysis}

\begin{table}[htbp]
\centering
\resizebox{\columnwidth}{!}{%
\begin{tabular}{llll}
\toprule
Layer Type & Cohen's d & 95\% CI & p-value \\
\midrule
down\_proj & 0.783 & [0.699, 0.869] & 5.89 $\times 10^{-26}$ \\
o\_proj & 0.772 & [0.686, 0.862] & 2.52 $\times 10^{-25}$ \\
v\_proj & 0.760 & [0.673, 0.851] & 1.23 $\times 10^{-24}$ \\
k\_proj & 0.759 & [0.660, 0.855] & 1.44 $\times 10^{-24}$ \\
up\_proj & 0.756 & [0.671, 0.844] & 1.95 $\times 10^{-24}$ \\
gate\_proj & 0.749 & [0.664, 0.836] & 4.96 $\times 10^{-24}$ \\
q\_proj & 0.745 & [0.649, 0.842] & 8.27 $\times 10^{-24}$ \\
\bottomrule
\end{tabular}
}
\end{table}

No layer type exceeds Cohen's d = 0.8. The range across layer types is narrow (0.745 to 0.783). The mean Cohen's d for attention layers (0.759) and MLP layers (0.763) are not significantly different (Mann-Whitney p = 1.0).

\begin{figure}[htbp]
\centering
\includegraphics[width=\columnwidth]{figures/cohens_d_by_layer_type.png}
\caption{Cohen's d by Layer Type}
\label{fig:cohens_d_by_layer_type}
\end{figure}
\textit{Figure 3: Cohen's d by layer type showing uniform clustering strength across all layer types.}

\subsubsection{P3 Control: Training Stochasticity}

\begin{table}[htbp]
\centering
\resizebox{\columnwidth}{!}{%
\begin{tabular}{ll}
\toprule
Metric & Value \\
\midrule
Within-task mean distance & 6.931 \\
Within-cluster mean distance & 7.786 \\
Ratio & 0.890 \\
Threshold & < 0.5 \\
Status & FAIL \\
\bottomrule
\end{tabular}
}
\end{table}

The P3 control criterion is not satisfied. The within-task variance (across different random seeds for the same task) is 89\% of the within-cluster variance (across different tasks in the same category). This indicates that training stochasticity contributes substantially to adapter geometry.

\subsubsection{Summary of Results}

\begin{table}[htbp]
\centering
\resizebox{\columnwidth}{!}{%
\begin{tabular}{llll}
\toprule
Research Question & Criterion & Result & Status \\
\midrule
RQ1: Existence & Cohen's d > 0.5 & d = 0.765 & PASS \\
RQ2: Mechanism & Spearman $\rho$ > 0.3 & $\rho$ = 0.389 & PASS \\
RQ3: Layer Analysis & Cohen's d > 0.8 for any layer & d = 0.783 (max) & FAIL \\
P3: Control & Ratio < 0.5 & Ratio = 0.89 & FAIL \\
\bottomrule
\end{tabular}
}
\end{table}
